\documentclass{article}
\usepackage{amsmath,amssymb}
\usepackage{xurl}
\usepackage[hidelinks]{hyperref}
% Shared commands for notes in docs/paper-gaps/.

\DeclareUnicodeCharacter{2080}{\ifmmode{}_{0}\else\textsubscript{0}\fi}
\DeclareUnicodeCharacter{2081}{\ifmmode{}_{1}\else\textsubscript{1}\fi}
\DeclareUnicodeCharacter{2082}{\ifmmode{}_{2}\else\textsubscript{2}\fi}
\DeclareUnicodeCharacter{2099}{\ifmmode{}_{n}\else\textsubscript{n}\fi}

\newcommand{\Lean}{\textsc{Lean}}
\ExplSyntaxOn
\NewDocumentCommand{\leanid}{m}{
  \ifmmode
    \textsf{\detokenize{#1}}
  \else
    \group_begin:
    \urlstyle{sf}
    \tl_set:Nn \l_tmpa_tl {#1}
    \tl_replace_all:Nnn \l_tmpa_tl {\_} {_}
    \exp_args:NV \path \l_tmpa_tl
    \group_end:
  \fi
}
\ExplSyntaxOff
\providecommand{\tr}{\operatorname{tr}}
\newcommand{\tnleanIssueBase}{https://github.com/LionSR/TNLean/issues/}
\newcommand{\tnleanPrBase}{https://github.com/LionSR/TNLean/pull/}
\newcommand{\tnleanDocBase}{https://sirui-lu.com/TNLean/docs/}
\newcommand{\ghissue}[1]{\href{\tnleanIssueBase#1}{GitHub issue \##1}}
\newcommand{\ghpr}[1]{\href{\tnleanPrBase#1}{PR \##1}}
\newcommand{\doclink}[1]{\href{\tnleanDocBase#1.html}{\path{#1}}}

% Dirac notation and the identity operator, as in blueprint/src/macros/common.tex.
% \providecommand keeps note-local definitions authoritative.
\usepackage{dsfont}
\providecommand{\ket}[1]{|#1\rangle}
\providecommand{\bra}[1]{\langle #1|}
\providecommand{\braket}[2]{\langle #1 | #2 \rangle}
\providecommand{\ketbra}[2]{|#1\rangle\!\langle #2|}
\providecommand{\Id}{\mathds{1}}
\providecommand{\one}{\mathds{1}}
\providecommand{\C}{\mathbb{C}}
\providecommand{\MN}[1]{M_{#1}(\C)}
\providecommand{\MD}{M_D(\C)}

% Verdict marker: \gapnote{<kind>}{<status>}. Kind is the policy
% classification (clarification, local-correction, scope-restriction,
% unfaithful, false-source, open-gap); status is open, wip (elimination
% underway), resolved, or historical. Machine-read by the site index;
% prints a verdict line.
\newcommand{\gapnote}[2]{\par\noindent\textbf{Verdict:} #1 (#2).\par}

\title{The Standard Form of a Matrix Product Unitary Is Taken in Canonical Form}
\author{}
\date{2026-11-16}
\begin{document}
\maketitle
\gapnote{local-correction}{resolved}

\section{The reading}

Propositions~\texttt{PropChiral} and \texttt{prop:local-char-transposition}
of Ref.~\cite{Cirac2017MPU} (source \path{Papers/1703.09188/paper_v2.tex},
lines 1214--1225 and 1266--1276) concern ``the MPU generated by a tensor in
SF'', and the sign lemma and Proposition~\texttt{prop:sigma-well-defined}
(lines 1595--1622) concern ``an MPU in SF represented by $u$ and $v$''. The
formalization reads ``a tensor in SF'' as the two-site block $\mathcal V_2$ of
a simple tensor $\mathcal V$ in canonical form~II, with any choice of the
gates $u,v$ and half factors in its standard form. Since $\mathcal V_2$ is
again in canonical form~II, these propositions then follow without further
hypotheses from the criteria for arbitrary gates in a canonical tensor.

\section{Evidence from the source}

The standard form is defined only through the canonical-form construction:
\begin{enumerate}
    \item Line 326 fixes the standing convention of Section~III: ``unless we
          specify the contrary, wlog we can assume that the tensor $\mathcal U$
          is in CF.'' Line 356 strengthens it to canonical form~II for
          $d^{-1/2}\mathcal U$ and all its blockings.
    \item The factors defining $u$ and $v$ are normalized by
          $X_1^\dagger(\Id\otimes\rho)X_1=\Id$ (lines 486--499), where $\rho$
          is the positive diagonal right fixed point of canonical form~II.
          Without canonical form this normalization, and hence the standard
          form (lines 603--622), is not defined.
    \item The Fundamental Theorem of MPU (lines 624--648), which the proof of
          \texttt{PropChiral} applies to standard forms, remarks that unitarity
          of $z$ ``follows from the fact that both tensors are in CFII''.
    \item Corollary~\texttt{coruvSF} (lines 870--884) starts from a path ``not
          necessarily in CF'' and puts in SF the blocked tensors of the CFII
          representatives $\hat{\mathcal U}(p)$, not the original tensors.
\end{enumerate}

The source's silence at Propositions~\texttt{PropChiral} and
\texttt{prop:local-char-transposition} is not a choice of the unrestricted
reading. Where the source intends tensors outside canonical form it says so
explicitly: ``not necessarily in CF'' at lines 710--712, 741, 747 and 871, and
``not necessarily in canonical form'' at line 1349. No such qualifier appears
at lines 1214--1288 or 1595--1631. The standing convention of line 326
therefore applies there.

\section{Consequence for the formalization}

A two-site standard form on an ambient tensor outside canonical form, for
example one padded by a zero virtual summand, is not a standard form in the
sense of Definition~\texttt{SF}. It is outside the scope of the source
propositions, so restricting to the canonical reading adds no hypothesis
absent from the source. The formal statements are
\leanid{MPOTensor.IsMPUCanonicalFormII.standardForm_even_timeReversal_iff}
and
\leanid{MPOTensor.IsMPUCanonicalFormII.standardForm_even_transposition_iff}.
The more general results for an arbitrary canonical ambient tensor keep their
own scope note, \cite{gap:mpu_canonical_form_full_support}.

The deformations in Proposition~\texttt{prop:sigma-well-defined} are paths of
tensors that need not be in canonical form, as in
Definition~\texttt{def:strictly-equivalent-symmetry} (line 1349). The sign at
each point is read from the standard form of the canonical representative,
which is how Corollary~\texttt{coruvSF} treats such paths.

The finite-chain hypotheses count sites of the two-spin standard-form chain;
see \cite{gap:mpu_standard_form_parity_gap}.

\bibliographystyle{alpha}
\bibliography{references}
\end{document}
